\documentclass[aps,prb,onecolumn,nofootinbib,superscriptaddress]{revtex4-2}
\usepackage{amsmath,amssymb,amsthm,mathtools,bm}
\usepackage{booktabs}
\usepackage[colorlinks=true,linkcolor=blue,citecolor=blue,urlcolor=blue]{hyperref}
\usepackage{microtype}

\newcommand{\src}[1]{\texttt{\detokenize{#1}}}
\newcommand{\E}{\mathbb{E}}
\newcommand{\Prob}{\mathbb{P}}
\newcommand{\dd}{\mathrm{d}}
\newcommand{\cA}{\mathcal{A}}
\newcommand{\cX}{\mathcal{X}}
\newcommand{\PF}{\mathfrak{P}}
\newcommand{\Koop}{\mathfrak{K}}

\begin{document}

\title{Supervisor 02 Notes: Core Dynamics Engines, Protocol Controls, and Production Provenance}
\author{Supervisor 02}
\date{June 11, 2026}

\begin{abstract}
This supervisor note synthesizes the assigned Worker 05--07 notes and an
independent source-level audit of the repository's core dynamics and protocol
surface.  The central conclusion is that the scientific dynamics should be
described through the canonical public drivers
\src{classA_U1FGTN.run_markov_circuit(...)} and
\src{classA_U1FGTN_gpu.run_markov_circuit(...)}, not through notebook-local
copies or private update helpers.  The physical perfect-correction protocol is a
Born-sampled adaptive Markov trajectory with deterministic correction after a
mismatch, whereas full post-selection is a deterministic forced-branch
trajectory and partial post-selection is an interpolation controlled by an
external Bernoulli mask.  The production layer has converged toward good
manifested GPU campaigns using streamed observers, but the repository still
contains stale Colab/OSG source copies, legacy scripts, incomplete
engine-internal canonical-entry metadata, and at least one serious OSG merge
cycle-label mismatch.  The Boss should preserve the protocol taxonomy and
canonical-entry architecture while scrutinizing source synchronization,
provenance fields, and legacy campaign paths before drawing final repository
claims.
\end{abstract}

\maketitle
\setcounter{tocdepth}{2}
\tableofcontents

\section{Source Scope}

I read \src{REPO_POLICY.md} first.  I then read only the three assigned worker
notes:
\begin{itemize}
\item \src{repo_synthesis/main_pass/worker_notes/worker_05_canonical_engines.tex};
\item \src{repo_synthesis/main_pass/worker_notes/worker_06_production_scripts_metadata.tex};
\item \src{repo_synthesis/main_pass/worker_notes/worker_07_protocol_taxonomy_controls.tex}.
\end{itemize}
I did not read other worker notes, supervisor notes, or boss notes.

The independent audit was read-only and source-level.  I inspected text/code
sources in \src{src/fgtn}, current non-excluded Colab source copies, selected
\src{scripts/} and \src{OSG/} production sources, \src{markov_transfer_operators},
the regularized Choi documentation, charge-fluctuation protocol documentation,
and \src{scripts/legacy/CI_Lindblad_DW.py}.  I did not edit files, execute notebooks, submit
jobs, run simulations, or load large binary arrays.  All inspected paths excluded
any path containing \src{haining}, \src{haoyu}, or \src{erroneous}, case
insensitively.

\section{Synthesis of Worker Claims}

\subsection{Canonical dynamics}

Worker 05 correctly identifies the repository policy's canonical production
surface: CPU production runs should call
\src{classA_U1FGTN.run_markov_circuit(...)} and GPU production runs should call
\src{classA_U1FGTN_gpu.run_markov_circuit(...)}.  The policy also forbids GPU
notebooks or scripts from reproducing the Markov cycle loop by directly calling
private helpers such as \src{_apply_grouped_site_updates(...)}.

The canonical CPU driver exposes the richest control surface: history/final-state
storage, sample parallelization, slab restriction, Choi diagnostics, and a
\src{physical_covariance_update} switch whose current default is the rank-one
resolvent update.  The canonical GPU driver is narrower by design: it is a
PyTorch Markov-circuit implementation with explicit device/dtype handling,
A100-oriented automatic batching, shard/manifest output, resume checks, snapshot
storage, streamed observers, Lyapunov observers, and Choi observers.

\subsection{Production scripts and metadata}

Worker 06's production-layer conclusion is also supported: recent GPU/OSG and
Colab campaigns generally use the canonical GPU public driver, keep covariance
histories out of memory when only observables are needed, and write campaign or
run metadata.  The strongest pattern is the observer-based campaign style:
\src{run_markov_circuit(..., save=False, return_data=False, cycle_observer=...)}
streams reduced observables while the campaign layer writes JSON summaries,
checkpoints, and manifests.

The production risks are operational rather than conceptual.  The modern OSG
slope-vs-cycle and slope-vs-system-size runners record
\src{canonical_dynamics_entry_point}, but the entropy-vs-time OSG runner calls
the canonical GPU method without storing that key in its block \src{config_json}.
The root \src{OSG/} path is stale: it has a placeholder container image,
hard-coded \src{queue 100} despite a generator currently producing four jobs,
and a wrapper call passing a \src{protocol} keyword that is not part of the
current canonical GPU signature.  The \src{OSG/osg_slope_vs_Ny} merge script
reconstructs \src{cycles = Ny // 2} and \src{fit_cycle = Ny // 2}, whereas its
generator and runner use \src{cycles = Ny} and \src{fit_cycle = Ny}; this is the
highest-impact source-level bug in the audited production layer.

\subsection{Protocol taxonomy}

Worker 07's taxonomy should be preserved nearly verbatim.  The physical
Born-sampled adaptive Markov process is a trajectory-resolved stochastic process
on covariance matrices.  Perfect correction does not remove the Born draw.  It
only makes the subsequent feedback/reset target deterministic once a mismatch is
observed.  Thus a perfect-correction observable is an ensemble average over
corrected Born trajectories, not a deterministic post-selected trajectory and not
an outcome-averaged transfer matrix.

Full post-selection is a forced desired-branch trajectory: upper-band modes are
forced empty and lower-band modes forced occupied.  Partial post-selection
interpolates by using an external Bernoulli mask to choose forced versus
Born-sampled site maps.  Transfer operators act on distributions or observables
of \(G\); Lyapunov data are quenched products along realized trajectories; the
regularized Choi object is a cumulative two-leg operator covariance for realized
branches; and the Lindblad object is a continuous-time mean/channel control, not
the sampled adaptive circuit.

\section{Independent Repo-Audit Adjustments}

\subsection{Source-copy synchronization is broader than Colab}

Worker 05 emphasizes stale Colab copies.  The independent audit confirms this
and extends it to bundled OSG copies.  The canonical GPU source
\src{src/fgtn/classA_U1FGTN_gpu.py} is byte-identical to only
\src{colab_regularized_choi_transfer_matrix/src/classA_U1FGTN_gpu.py}.  Five
other Colab GPU copies and two OSG GPU copies share the older checksum
\src{3196fa36...c2216} and have 2831 lines, while the canonical source has
checksum \src{6b90e1e0...a112c5} and 2881 lines.  The older copies also expose
older Choi metadata, including \src{regularized_rank_one_basis_independent_v1}
rather than the canonical \src{regularized_resolvent_rank_one_v2}.  Therefore
the Boss should not treat all non-excluded \src{classA_U1FGTN_gpu.py} files as
synchronized.

The CPU source-copy issue is narrower but real: the canonical CPU source has
5308 lines and checksum \src{f0dcd7ab...c5f6}, whereas the Colab charge
fluctuation CPU copy has 4779 lines and checksum \src{6816fa01...897}.  Any
claim about CPU algorithmic details should be anchored to \src{src/fgtn}, unless
a result explicitly records which copied source produced it.

\subsection{Canonical-entry provenance is split between engines and campaigns}

Worker 05 notes that the canonical CPU/GPU \src{run_markov_circuit} configs do
not visibly include \src{canonical_dynamics_entry_point}.  I confirm that the
CPU \src{run_markov_circuit} \src{_run_config} records physics, protocol, Choi,
and slab fields but not that explicit key, and the GPU \src{run_config} similarly
omits it.  The CPU \src{run_markov_channel} routine does record
\src{canonical_dynamics_entry_point}, so the provenance pattern exists in the
codebase.

The adjustment is that some campaign layers add the canonical-entry key outside
the engine.  For example, current small-system GPU run summaries and Lyapunov
campaign JSON files record \src{classA_U1FGTN_gpu.run_markov_circuit}, and charge
fluctuation CPU run summaries can record \src{classA_U1FGTN.run_markov_circuit}.
Thus the provenance gap is not universal at the output layer.  The precise
statement should be: engine-internal Markov-circuit manifests/configs are not
self-sufficient with respect to the policy-mandated canonical-entry key, but many
modern campaign wrappers repair this at the run-summary or campaign-manifest
layer.

\subsection{Partial post-selection is now part of the canonical control surface}

Worker 07 describes the partial-post-selection fork.  The independent audit finds
the same intermediate \src{postselect_probability} logic in the canonical CPU and
GPU engines as well: \src{postselect=True} is normalized to probability one,
probabilities outside \([0,1]\) are rejected, full post-selection forces
\src{samples=1}, and intermediate probabilities split site batches by a random
mask.  The partial-post-selection campaign tree remains a useful concrete
application, but the protocol knob itself is not only fork-local in the current
canonical source.

\subsection{Modern production style is real, but legacy scripts remain visible}

The modern production path is consistent: OSG and Colab runners stream
observables through \src{cycle_observer}, \src{site_observer},
\src{lyapunov_observer}, or Choi observers rather than saving full histories by
default.  The older \src{scripts/} layer is mixed.  Some scripts correctly call
the CPU canonical \src{run_markov_circuit}; others still call
\src{run_adaptive_circuit} or save only figures/raw histories without a separate
machine-readable manifest.  These older scripts should be treated as exploratory
or legacy unless a specific run summary and source checksum are present.

\section{Key Evidence Table}

\begin{table}[h]
\caption{Source-level evidence from the assigned workers and independent audit.}
\begin{ruledtabular}
\begin{tabular}{p{0.30\linewidth}p{0.16\linewidth}p{0.46\linewidth}}
Source & Lines & Evidence and implication \\
\hline
\src{REPO_POLICY.md} & 5--24, 42--45 &
Defines \src{classA_U1FGTN.run_markov_circuit(...)} and
\src{classA_U1FGTN_gpu.run_markov_circuit(...)} as canonical production entry
points; forbids direct GPU use of private update helpers; requires production
metadata to record the canonical entry point; requires non-excluded GPU source
copies to remain synchronized. \\
\src{src/fgtn/classA_U1FGTN.py} & 1488--1500, 1666--1714 &
CPU Markov site update samples Born outcomes before feedback.  With
\src{perfect_correction=True}, the feedback target is forced to pump out upper
mismatches or pump in lower mismatches; the measurement outcome itself is not
post-selected. \\
\src{src/fgtn/classA_U1FGTN.py} & 2438--2491 &
Canonical CPU \src{run_markov_circuit} signature exposes history/final storage,
postselection probability, perfect correction, sample parallelization, slab
restriction, Choi observers, rank-one/dense physical update selection, and cycle
observers. \\
\src{src/fgtn/classA_U1FGTN.py} & 2513--2621 &
CPU driver validates \src{postselect_probability}, normalizes full
post-selection, records \src{physical_covariance_update},
\src{choi_formula}, slab indices, and exterior preparation, but does not visibly
record \src{canonical_dynamics_entry_point} in this Markov-circuit run config. \\
\src{src/fgtn/classA_U1FGTN.py} & 3237--3259 &
Separate CPU \src{run_markov_channel} config records
\src{canonical_dynamics_entry_point}, showing the provenance convention exists
elsewhere in the source. \\
\src{src/fgtn/classA_U1FGTN_gpu.py} & 15--22, 76--125 &
GPU class is explicitly a narrow Markov-circuit implementation and guards
PyTorch/CUDA device resolution. \\
\src{src/fgtn/classA_U1FGTN_gpu.py} & 228--292 &
GPU batch-size estimator targets an A100 40 GB budget, adjusts for detected free
memory, storage mode, dtype/backend, snapshots, and Choi tracking. \\
\src{src/fgtn/classA_U1FGTN_gpu.py} & 1913--2018 &
Private grouped-site helper dispatches full post-selection, pure Born Markov
updates, and intermediate partial post-selection by a random mask.  This is
engine-internal machinery, not a production entry point. \\
\src{src/fgtn/classA_U1FGTN_gpu.py} & 2231--2430 &
Canonical GPU \src{run_markov_circuit} validates resume, Choi observer
compatibility, postselection probability, snapshots, batching, and records rich
run configuration including \src{physical_covariance_update} and
\src{choi_formula}, but not an explicit canonical-entry key. \\
\src{src/fgtn/classA_U1FGTN_gpu.py} & 2586--2626 &
The public GPU driver owns the cycle loop and calls the private grouped-site
helper internally while appending snapshots/history according to storage mode. \\
Checksum audit & N/A &
Canonical GPU source and the regularized-Choi Colab copy share checksum
\src{6b90e1e0...a112c5}; five other Colab GPU copies plus two OSG GPU copies
share older checksum \src{3196fa36...c2216}. \\
Line-count audit & N/A &
Canonical GPU source has 2881 lines; older Colab/OSG GPU copies have 2831 lines.
Canonical CPU source has 5308 lines; the charge-fluctuation Colab CPU copy has
4779 lines. \\
\src{OSG/scripts/run_slope_vs_cycle_block.py} & 65--88, 183--197 &
Modern OSG slope-vs-cycle runner records
\src{canonical_dynamics_entry_point} and uses canonical GPU
\src{run_markov_circuit} with \src{save=False} and a streamed cycle observer. \\
\src{OSG/osg_slope_vs_Ny/generate_params.py} & 41--57 &
System-size OSG parameter generation writes \src{cycles=Ny} and
\src{fit_cycle=Ny}. \\
\src{OSG/osg_slope_vs_Ny/run_slope_vs_system_size.py} & 55--83, 124--132 &
System-size OSG runner records canonical-entry metadata and calls canonical GPU
\src{run_markov_circuit}. \\
\src{OSG/osg_slope_vs_Ny/merge_slope_vs_system_size.py} & 54--80 &
Merge script reconstructs \src{cycles=Ny//2} and \src{fit_cycle=Ny//2}, in
conflict with the generator/runner convention. \\
\src{OSG/osg_entropy_vs_time_maxmix/run_entropy_vs_time.py} & 106--137 &
Entropy-vs-time runner calls canonical GPU dynamics, but the saved
\src{config_json} contains no explicit canonical-entry key. \\
\src{OSG/osg_entropy_vs_time_maxmix/merge_entropy_vs_time.py} & 31--49 &
Merge checks cycle-array equality but uses \src{count[0]} as the reported sample
count, so uneven count completeness across cycles is not fully asserted. \\
\src{OSG/generate_params.py}, \src{OSG/batch_gpu.submit}, \src{OSG/run.sh} &
48--58; 1--28; 38--46 &
Legacy root OSG path generates four jobs but queues 100, uses a placeholder
container image, and calls \src{run_markov_circuit} with a \src{protocol} keyword
that is not in the current canonical GPU signature. \\
\src{scripts/generate_markov_history_samples.py} & 30--44, 99--113 &
Representative CPU script sets CPU/thread controls and calls canonical CPU
\src{run_markov_circuit}. \\
\src{scripts/generate_G_from_existing_G0.py} & 105--133 &
Older generation script calls \src{run_adaptive_circuit} and saves raw history
plus source snapshot metadata, not canonical Markov-circuit provenance. \\
\src{scripts/analyze_covariance_protocol_histories_characterization.py} &
36--48, 325--340 &
Downstream analysis expects run summaries to contain
\src{classA_U1FGTN_gpu.run_markov_circuit}, \src{store_mode=history}, one shard,
and no snapshots, demonstrating campaign-layer provenance validation. \\
\src{markov_transfer_operators/markov_transfer_operators_note.tex} &
133--215, 247--299 &
Formalizes the covariance dynamics as a Markov chain with kernel composition and
dual Koopman/Perron-Frobenius transfer actions. \\
\src{markov_transfer_operators/markov_transfer_operators_note.tex} &
370--461 &
Lyapunov exponents are quenched singular-value growth rates of chronological
branchwise tangent products, not spectra of outcome-averaged matrices. \\
\src{colab_regularized_choi_transfer_matrix/docs/regularized_choi_covariance_contractions.tex} &
63--103, 105--164 &
Choi object is a cumulative two-leg Gaussian operator covariance for realized
branches and is explicitly distinguished from production state-covariance
dynamics. \\
\src{scripts/legacy/CI_Lindblad_DW.py} & 516--601, 727--858 &
Lindblad control evolves a single-layer covariance by RK4 through gain/loss and
optional decoherence and constructs an affine vectorized mean-channel generator. \\
\src{colab_charge_fluctuations/docs/charge_fluctuations_notes.tex} &
274--299 &
Protocol notes describe perfect correction as a multi-sample ensemble average
and post-selection as a single fixed-pattern trajectory. \\
\end{tabular}
\end{ruledtabular}
\end{table}

\section{Contradictions and Caveats}

There is no contradiction between ``perfect correction'' and Born sampling once
the word ``perfect'' is read correctly.  Perfect correction means perfect
classical correction/reset after a mismatch, not perfect selection of the
measurement outcome.  Any Boss-level prose that calls perfect correction a
post-selected trajectory, a deterministic covariance map, or an averaged transfer
matrix should be revised.

There is an apparent provenance tension.  The repository policy requires every
production output to record the canonical dynamics entry point.  The canonical
CPU/GPU Markov-circuit engine configs do not visibly record that field, but many
modern campaign summaries do.  The Boss should distinguish engine-level metadata
from campaign-layer metadata and should prefer outputs whose run summaries or
manifests explicitly include both the canonical entry point and the source copy
or checksum.

There is also an apparent source-identity tension.  A Colab or OSG campaign may
record \src{classA_U1FGTN_gpu.run_markov_circuit} while using a stale copied
\src{classA_U1FGTN_gpu.py}.  The canonical-entry string proves use of the public
method name, but it does not by itself prove byte identity with \src{src/fgtn}.
For high-stakes comparisons, source checksums should be preserved.

The production audit is source-only.  I did not execute notebooks, run scripts,
submit OSG jobs, compare numerical outputs, or load covariance arrays.  Runtime
failure claims are therefore based on visible source inconsistencies, not
observed executions.

Line numbers refer to the current working tree inspected for this note.  They
may drift if files are edited later.

\section{Confidence}

Confidence is high for the canonical-entry identification, the protocol
taxonomy, and the source-copy synchronization findings, because these follow
directly from policy text, engine signatures, branch-control code, line counts,
and checksums.  Confidence is high that the OSG
\src{cycles=Ny} versus \src{cycles=Ny//2} mismatch is a real source-level defect.
Confidence is moderate for production-status classifications such as ``modern''
or ``legacy,'' since those are inferred from directory organization, metadata
style, and runner freshness rather than an external registry.  Confidence is low
for any numerical-performance or physics-output claim, because no simulations or
notebooks were executed.

\section{What the Boss Should Preserve or Scrutinize}

\subsection{Preserve}

The Boss should preserve the canonical public-entry narrative:
\[
  \src{classA_U1FGTN.run_markov_circuit(...)}
  \quad\text{and}\quad
  \src{classA_U1FGTN_gpu.run_markov_circuit(...)}
\]
are the authoritative CPU/GPU adaptive Markov-circuit surfaces.  Private helpers
belong below these methods.

The Boss should preserve the protocol taxonomy: Born Markov trajectory,
perfect-correction ensemble, full post-selection, partial post-selection,
distributional transfer operators, Choi operator-transfer objects, and Lindblad
mean-channel controls are distinct objects.

The Boss should preserve the modern campaign pattern: streamed observers,
\src{save=False} where full histories are unnecessary, shard/checkpoint
manifests, explicit protocol fields, and campaign-level summaries.

\subsection{Scrutinize}

The Boss should scrutinize all noncanonical copied
\src{classA_U1FGTN_gpu.py} and \src{classA_U1FGTN.py} files, especially stale
Colab and OSG copies, before using them as evidence for current dynamics.

The Boss should require explicit canonical-entry metadata and source identity
for production outputs.  A method-name string alone is weaker than a method-name
string plus source checksum or copied-source manifest.

The Boss should treat the root \src{OSG/} path as stale unless repaired, and
should specifically flag the \src{OSG/osg_slope_vs_Ny} merge cycle mismatch.

The Boss should avoid conflating Choi, Lindblad, transfer-operator, and sampled
trajectory language.  These controls are useful cross-checks and theoretical
objects, but they are not interchangeable with the production state-covariance
Markov process.

\section{Cross-Supervisor Addendum}

\subsection{Cross-note agreements}

All five supervisor notes agree on the central separation of objects: the
physical adaptive circuit is a trajectory-resolved Born process, perfect
correction is an ensemble of Born trajectories with deterministic feedback after
mismatch, and post-selection is a deterministic or one-sample control rather
than an equivalent ensemble.  This reinforces the Supervisor 02 taxonomy and the
ensemble distinction
\[
\overline{O[G_\xi]} \neq O[\overline{G_\xi}] \neq O[G_{\rm ps}] .
\]

The notes also agree that the strongest evidence spine is not raw cache data but
the structured campaign layer: canonical-entry metadata, run summaries,
manifests, run-index CSVs, explicit sample counts, and derived-analysis
manifests.  Supervisor 05 independently strengthens the Supervisor 02 warning
that a method-name string is insufficient unless tied to source identity,
campaign metadata, and sample-count semantics.

The target-state and trajectory notes are compatible with the dynamics note:
OW/domain-wall target diagnostics provide the theoretical benchmark, while
perfect-correction trajectories supply the physical ensemble evidence.  Entropy
scaling near \(1/3\) is the leading critical diagnostic; correlator exponents,
partial-postselection sweeps, and finite-size transition claims require stronger
qualification.

The transfer-stability note agrees with Supervisor 02 that Choi, Lindblad,
Markov-transfer, and sampled covariance dynamics are distinct.  It adds that
Lyapunov and Choi transfer gaps are complementary diagnostics of a near-marginal
domain-wall sector, not one mathematically identical spectrum.

\subsection{Corrections and refinements to Supervisor 02}

There is no hard cross-note contradiction to the Supervisor 02 protocol
taxonomy.  The main correction is emphasis: perfect correction should not be
described as fixed-charge projection.  Supervisor 03 shows strong suppression of
\(\operatorname{Tr} C(1-C)\), but centered charge wander remains small and
nonzero.

Supervisor 02's provenance caution should be sharpened with Supervisor 05's
figure/data caveats.  The Boss should not only require canonical-entry and source
checksum metadata, but also check actual versus requested samples, skipped-input
records, figure-script source tables, and whether a plotted comparison is backed
by matching protocol data.

The partial-postselection control should remain in the taxonomy, but Supervisor
03 shows that the inspected sweep is finite-geometry protocol evidence, not a
phase diagram.  The \(p=1\) endpoint is effectively a deterministic
postselected branch.

The Choi discussion in Supervisor 02 should inherit Supervisor 04's stricter
language: GPU Choi evidence is early-cycle and censored before final time, while
CPU Choi evidence reaches final time but has endpoint/finite-mode caveats.  Do
not collapse these into a single final ``transfer gap.''

The target-state context from Supervisor 01 adds one cross-folder convention
warning absent from Supervisor 02: \(\alpha_1,\alpha_2\) are not used uniformly
across theory folders.  Boss-level prose should say ``topological mass'' and
``trivial mass'' unless the local convention is stated.

\subsection{Boss priorities}

\begin{enumerate}
\item Preserve the canonical dynamics and ensemble taxonomy: CPU/GPU
\src{run_markov_circuit} entry points, Born perfect-correction ensemble,
postselection control, partial-postselection interpolation, and distinct
Choi/Lyapunov/Lindblad/Markov-transfer objects.

\item Use structured manifests and run summaries as the data spine.  Require
canonical entry point, source identity or checksum, actual sample count, storage
mode, completion status, and source-to-derived-product links.

\item Treat entropy scaling and purification as the strongest physics evidence;
scrutinize correlator exponents, postselection comparisons, partial-postselection
phase language, and any exact-preparation claim.

\item Repair or quarantine stale operational paths: unsynchronized Colab/OSG
source copies, root \src{OSG/}, the \src{OSG/osg_slope_vs_Ny} cycle mismatch,
legacy \src{cache/G_history_samples}, and hard-coded or mismatched summary-figure
inputs.

\item Keep target-state and stability claims precise: OW one-shell captures
chiral entanglement better than full Chern quantization, DW truncation is an
active boundary condition, and Lyapunov/Choi gaps support a qualitative
near-marginal domain-wall sector but not a resolved thermodynamic exponent.
\end{enumerate}

\end{document}
